\documentclass[11pt]{article}
\usepackage[margin=1in]{geometry}
\usepackage{times}

\begin{document}

\begin{center}
{\Large\bfseries Ultra-Lightweight Dense and Recurrent 2D Neural Networks for Edge and Mobile Image Recognition}

\vspace{0.5em}
Babak Shiri$^{1}$ and Zahra Alijani$^{2}$

\vspace{0.3em}
\small
$^{1}$Key Laboratory of Numerical Simulation of Sichuan Province, School of Mathematics and Big Data, Neijiang Normal University, Neijiang, 641100, Sichuan, China\\
$^{2}$Institute for Research and Applications of Fuzzy Modeling, University of Ostrava, Ostrava, 701 03, Czech Republic
\end{center}

\vspace{1em}
\noindent\textbf{Abstract}

\noindent Ultra-lightweight image recognition is studied using neural networks that keep an image's matrix structure instead of flattening it into a vector, processing it directly through paired left and right linear transformations. Building on this bilinear formulation -- related to prior matrix neural network and matrix-RNN work -- both a dense two-layer architecture and a patch-based recurrent variant are developed and systematically compared across a three-tier parameter budget (super-tiny, tiny, ultra-light). Both variants are evaluated on the MNIST and Fashion-MNIST datasets, reaching competitive classification accuracy with extremely low parameter counts and fast training. A real-time Android application is implemented for handwritten digit recognition, demonstrating practical deployability on resource-constrained edge devices.

\vspace{1em}
\noindent\textbf{Keywords:} 2D neural networks, recurrent neural network, lightweight deep learning, ultra-light networks, MNIST classification, mobile implementation, edge AI

\end{document}
